\section{Algorithm Description}

\subsection{Formalising the search problem}

Bloxorz is fully observable and deterministic: the whole board is visible and each action
leads to exactly one successor. We therefore model it as a classical search problem. All four
solvers share the same problem definition and never re-implement the game rules; they only call
\texttt{successors(level, state)} and \texttt{is\_goal(level, state)} from
\texttt{engine/transitions.py}. A rule therefore lives in one place and affects every algorithm
in the same way, which keeps the comparison fair.

\subsubsection{State representation}

The static stage (tiles, goal, switch definitions) is stored in a \texttt{Level} object and is
\emph{not} part of the state. A state is an immutable \texttt{GameState}
(Table~\ref{tab:state}).

\begin{table}[H]
\centering
\caption{Components of a search state}
\label{tab:state}
\begin{tabularx}{\linewidth}{llX}
\toprule
\textbf{Field} & \textbf{Type} & \textbf{Meaning} \\
\midrule
\texttt{block.orientation} & \texttt{S}, \texttt{H}, \texttt{V} & Standing, lying horizontally, lying vertically \\
\texttt{block.row}, \texttt{block.col} & int & \texttt{S}: occupied cell; \texttt{H}: left cell; \texttt{V}: top cell \\
\texttt{bridge\_mask} & int (bitmask) & Bit $i$ is 1 when bridge group $i$ is open \\
\texttt{cubes} & \texttt{None} or two cells & Positions of the two cubes after a split \\
\texttt{active\_cube} & 0 or 1 & The cube that the next move controls \\
\bottomrule
\end{tabularx}
\end{table}

\begin{itemize}[leftmargin=*]
\item \textbf{Orientation} is part of the state because it decides which cells are occupied,
whether the block may rest on a fragile tile, whether it can win, and how each roll transforms it.
\item \textbf{Bridge configuration} is part of the state, as required: the same block placement
with different open/closed bridges allows different moves. A bitmask keeps the state small and
fast to hash; with $B$ bridge groups there are at most $2^B$ configurations.
\item \textbf{Split mode} needs \texttt{cubes} and \texttt{active\_cube}; the active cube decides
which cube the next action moves.
\item \texttt{GameState} is a frozen dataclass, hence hashable, so whole states are used as
dictionary keys and set members to detect repeated states.
\end{itemize}

For a stage with $N$ floor cells, an unsplit block has at most about $3N$ placements, multiplied
by $2^B$ bridge configurations; split mode adds at most $2N(N-1)$ cube configurations. Our test
stages have 37 to 90 floor cells and at most two bridge groups, so only hundreds to a few thousand
states are reachable; the \emph{Test All} stage (432 floor cells, six bridge groups, a split
switch) is much larger.

\subsubsection{Initial state and goal test}

The initial state is the state passed to the solver: the GUI passes the current state, which may
be mid-level, split, or have bridges toggled; the experiments use the stage's start state (block
standing on \texttt{S}, bridges as defined in the JSON file). The goal test holds only when the
block is whole, \emph{standing}, and on the goal cell. A block lying across the goal and a single
cube do not win.

\subsubsection{Actions and transition model}

The actions are \texttt{UP}, \texttt{DOWN}, \texttt{LEFT}, \texttt{RIGHT} and, only while the
block is split, \texttt{SWAP}. \texttt{transition(level, state, action)} returns the next state
or \texttt{None} for an illegal move, in five steps.

\textbf{(1) Orientation table.} Table~\ref{tab:roll} lists the twelve deterministic
(orientation, direction) cases. A standing block tips over in the rolling direction, a block
lying along the rolling direction stands up, and a block lying across it rolls sideways.

\begin{table}[H]
\centering
\caption{Orientation transition table, block at $(r,c)$}
\label{tab:roll}
\begin{tabular}{lllll}
\toprule
\textbf{From} & \textbf{UP} & \textbf{DOWN} & \textbf{LEFT} & \textbf{RIGHT} \\
\midrule
\texttt{S} & \texttt{V} $(r-2,c)$ & \texttt{V} $(r+1,c)$ & \texttt{H} $(r,c-2)$ & \texttt{H} $(r,c+1)$ \\
\texttt{H} & \texttt{H} $(r-1,c)$ & \texttt{H} $(r+1,c)$ & \texttt{S} $(r,c-1)$ & \texttt{S} $(r,c+2)$ \\
\texttt{V} & \texttt{S} $(r-1,c)$ & \texttt{S} $(r+2,c)$ & \texttt{V} $(r,c-1)$ & \texttt{V} $(r,c+1)$ \\
\bottomrule
\end{tabular}
\end{table}

\textbf{(2) Legality.} Every occupied cell must be supported (not void, not off the board, not a
closed bridge), and a standing block must not be on a fragile tile.

\textbf{(3) Switches} fire only on a fresh press (pressed after the move but not before it).
A soft switch fires when any part of the block, or a cube, covers it; a heavy switch only when
the block stands on it. A \emph{toggle} effect flips the bridge bits (XOR); a \emph{permanent}
effect sets them to a fixed state.

\textbf{(4) Split.} A block that arrives standing on a split switch is teleported into two cubes
at predefined cells (optionally changing bridges). Each move then moves the active cube by one
cell; cubes can press soft switches only.

\textbf{(5) Merge.} When the two cubes become adjacent they merge into an \texttt{H} or
\texttt{V} block, if that placement is legal.

Illegal moves never appear in \texttt{successors()}, so no solver generates a falling state.

\subsubsection{Cost function}

BFS and DFS ignore costs. UCS and A* share the risk-aware cost \texttt{move\_cost} defined in
\texttt{solvers/common.py} (Table~\ref{tab:cost}).

\begin{table}[H]
\centering
\caption{Step cost used by UCS and A*}
\label{tab:cost}
\begin{tabular}{lc}
\toprule
\textbf{Transition} & \textbf{Cost} \\
\midrule
\texttt{SWAP} & 0 \\
Each roll of the block or step of a cube & 1 \\
Additional penalty if any occupied cell ends on a fragile tile & $+4$ \\
\bottomrule
\end{tabular}
\end{table}

\begin{itemize}[leftmargin=*]
\item \textbf{Why not uniform.} With a cost of 1 per move, UCS reduces to BFS. We make risky
moves expensive instead.
\item \textbf{Fragile penalty.} A fragile tile breaks if the block ends up standing on it, so a
path that keeps landing on fragile tiles is one wrong orientation away from failure. A penalty of
4 means the solver accepts a detour of up to three extra rolls to avoid one fragile landing. It is
large enough to change the path when a detour exists and small enough that unavoidable fragile
tiles are still crossed.
\item \textbf{SWAP is free} because it moves nothing; it only changes which cube is controlled.
This also matches the GUI's Moves counter, which does not count \texttt{SWAP}.
\item \textbf{Consequences.} All costs are non-negative, which UCS and A* need for optimality,
and every non-\texttt{SWAP} move costs at least 1, which the A* heuristic relies on
(Section~\ref{sec:heuristic}). UCS and A* minimise risk-weighted cost, not the number of moves;
the number of moves is still reported separately.
\end{itemize}

\subsubsection{Avoiding repeated states}

All four algorithms are graph searches and always compare the \emph{whole} \texttt{GameState},
never just coordinates (Table~\ref{tab:visited}). The state space is finite and each state is
expanded at most once, so no algorithm can loop. Two consecutive zero-cost \texttt{SWAP}s return
to a state that is already closed, so they cannot loop either. A dictionary
\texttt{parents[state] = (previous\_state, action)} rebuilds the solution by walking back from the
goal and reversing the actions.

\begin{table}[H]
\centering
\caption{Repeated-state handling}
\label{tab:visited}
\begin{tabular}{lll}
\toprule
\textbf{Algorithm} & \textbf{Structure} & \textbf{State marked when} \\
\midrule
BFS & \texttt{parents} dictionary & generated \\
DFS & \texttt{parents} dictionary & popped from the stack \\
UCS, A* & \texttt{best} cost dictionary $+$ \texttt{closed} set & popped from the priority queue \\
\bottomrule
\end{tabular}
\end{table}

\subsubsection{Measurements}

Every solver returns a \texttt{SearchResult} with the four required measurements:
\textbf{search time} (\texttt{time.perf\_counter()} around the search loop),
\textbf{peak memory} (the peak of Python allocations during the search, from
\texttt{tracemalloc}), \textbf{expanded nodes} (the number of states whose successors were
generated) and \textbf{solution length} (the number of moves, \texttt{SWAP} excluded).

% -------------------------------------------------------------------------
\subsection{Breadth-First Search (BFS)}

BFS expands states level by level: first every state one action away from the start, then every
state two actions away, and so on. The frontier is a FIFO queue (\texttt{collections.deque}).

\begin{algorithm}[H]
\caption{BFS (\texttt{solvers/bfs.py})}
\label{alg:bfs}
\begin{algorithmic}[1]
\If{\Call{IsGoal}{$start$}} \Return empty path \EndIf
\State $parents \gets \{start: \textsc{None}\}$; $frontier \gets$ queue $[start]$
\While{$frontier$ is not empty}
  \State $s \gets frontier.\textsc{PopLeft}()$; $expanded \gets expanded + 1$
  \For{each $(a, s')$ in \Call{Successors}{$s$}}
    \If{$s' \in parents$} \textbf{continue} \EndIf
    \State $parents[s'] \gets (s, a)$
    \If{\Call{IsGoal}{$s'$}} \Return \Call{Rebuild}{$parents, s'$} \Comment{goal test on generation}\EndIf
    \State $frontier.\textsc{Append}(s')$
  \EndFor
\EndWhile
\State \Return no solution
\end{algorithmic}
\end{algorithm}

\begin{itemize}[leftmargin=*]
\item \texttt{parents} doubles as the visited set, so every state enters the queue once.
\item The goal test is applied when a state is \emph{generated}. With BFS the first generated goal
is always at minimum depth, so stopping early keeps optimality and saves a whole layer of
expansions.
\item \textbf{Complete}: the state space is finite. \textbf{Optimal in the number of actions},
\texttt{SWAP} included. It does not guarantee the fewest displayed moves (10 moves $+$ 3
\texttt{SWAP} $=$ 13 actions loses to 11 moves $+$ 0 \texttt{SWAP}) and it ignores fragile tiles.
\item Time and memory $O(b^d)$, where $b \le 4$ ($5$ when split) is the branching factor and $d$
the solution depth; in practice bounded by the number of reachable states.
\end{itemize}

% -------------------------------------------------------------------------
\subsection{Depth-First Search (DFS)}

We implement \textbf{graph-search DFS} (not IDS). DFS always follows the newest branch until it
gets stuck, then backtracks. The frontier is a LIFO stack.

\begin{algorithm}[H]
\caption{DFS (\texttt{solvers/dfs.py})}
\label{alg:dfs}
\begin{algorithmic}[1]
\State $parents \gets \{\}$; $stack \gets [(start, \textsc{None})]$
\While{$stack$ is not empty}
  \State $(s, link) \gets stack.\textsc{Pop}()$
  \If{$s \in parents$} \textbf{continue} \EndIf
  \State $parents[s] \gets link$ \Comment{marked when popped}
  \If{\Call{IsGoal}{$s$}} \Return \Call{Rebuild}{$parents, s$} \EndIf
  \State $expanded \gets expanded + 1$
  \For{each $(a, s')$ in \textsc{Reversed}(\Call{Successors}{$s$})}
    \If{$s' \notin parents$} $stack.\textsc{Push}((s', (s, a)))$ \EndIf
  \EndFor
\EndWhile
\State \Return no solution
\end{algorithmic}
\end{algorithm}

\begin{itemize}[leftmargin=*]
\item An explicit stack replaces recursion, so long paths (1363 moves on Test All) cannot hit
Python's recursion limit.
\item A state is marked when \emph{popped}, not when pushed, so the recorded parent link is the
route DFS actually followed to it.
\item Successors are pushed in reverse so that the engine's first action (\texttt{UP}, then
\texttt{DOWN}, \texttt{LEFT}, \texttt{RIGHT}, \texttt{SWAP}) is explored first.
\item \textbf{Complete} on a finite space thanks to the visited set, but \textbf{not optimal}: it
returns the first solution it meets. The frontier stays small ($O(b\,m)$ for maximum depth $m$),
although the visited set can grow to the size of the state space.
\end{itemize}

% -------------------------------------------------------------------------
\subsection{Uniform-Cost Search (UCS)}

UCS expands the state with the smallest path cost $g(n)$ first. The frontier is a binary heap
(\texttt{heapq}) keyed on $g$.

\begin{algorithm}[H]
\caption{UCS (\texttt{solvers/ucs.py})}
\label{alg:ucs}
\begin{algorithmic}[1]
\State $best \gets \{start: 0\}$; $parents \gets \{start: \textsc{None}\}$; $closed \gets \emptyset$
\State $heap \gets [(0, 0, tie, start)]$ \Comment{(cost, number of actions, insertion order, state)}
\While{$heap$ is not empty}
  \State $(g, depth, \_, s) \gets heap.\textsc{PopMin}()$
  \If{$s \in closed$} \textbf{continue} \Comment{stale entry}\EndIf
  \If{\Call{IsGoal}{$s$}} \Return \Call{Rebuild}{$parents, s$} \Comment{goal test on pop}\EndIf
  \State $closed \gets closed \cup \{s\}$; $expanded \gets expanded + 1$
  \For{each $(a, s')$ in \Call{Successors}{$s$}}
    \If{$s' \in closed$} \textbf{continue} \EndIf
    \State $g' \gets g + \Call{MoveCost}{s, s', a}$
    \If{$s' \notin best$ \textbf{or} $g' < best[s']$}
      \State $best[s'] \gets g'$; $parents[s'] \gets (s, a)$
      \State $heap.\textsc{Push}((g', depth + 1, tie, s'))$
    \EndIf
  \EndFor
\EndWhile
\State \Return no solution
\end{algorithmic}
\end{algorithm}

\begin{itemize}[leftmargin=*]
\item \textbf{Goal test on pop.} With non-uniform costs, the goal may first be generated through
an expensive route; only when it is popped is its $g$ guaranteed minimal.
\item \textbf{Lazy deletion.} \texttt{heapq} has no decrease-key, so a cheaper route pushes a new
entry and the old one is skipped when popped (it is already in \texttt{closed}).
\item \textbf{Tie-breaking.} Among equal costs, fewer actions are preferred; the insertion counter
prevents Python from comparing two \texttt{GameState} objects.
\item \textbf{Complete and optimal} with respect to \texttt{move\_cost}, because all step costs are
non-negative. Time $O(N \log N)$ and memory $O(N)$, where $N$ is the number of states cheaper than
the optimal cost $C^*$. UCS grows in cost contours, so when the solution must cross a fragile tile
it also expands states deeper than the solution depth.
\end{itemize}

% -------------------------------------------------------------------------
\subsection{A* Search}

A* uses the same skeleton and the same cost $g$ as UCS, but orders the heap by
\begin{equation}
f(n) = g(n) + h(n),
\end{equation}
where $h(n)$ estimates the remaining cost. The heap key is $(f, h, \text{insertion order})$, so
among equal $f$ the state closer to the goal is preferred. Because $h$ is consistent
(Section~\ref{sec:admissible}), each state can be closed on its first expansion without reopening.

\subsubsection{Heuristic design}
\label{sec:heuristic}

\textbf{Idea: a lower bound on moving the block's centre.} Ignore terrain, fragile tiles, bridges
and orientation legality, and ask only how many rolls are needed to move the block's centre to the
goal's centre. Coordinates are doubled (half-cell units) so that centres are integers
(Table~\ref{tab:centre}); the goal is a standing block on $G=(g_r, g_c)$, with centre
$(2g_r, 2g_c)$.

\begin{table}[H]
\centering
\caption{Block centre in doubled coordinates}
\label{tab:centre}
\begin{tabular}{ll}
\toprule
\textbf{Shape} & \textbf{Centre} \\
\midrule
\texttt{S} at $(r,c)$ & $(2r, 2c)$ \\
\texttt{H} at $(r,c)$ & $(2r, 2c+1)$ \\
\texttt{V} at $(r,c)$ & $(2r+1, 2c)$ \\
Two cubes at $(r_0,c_0), (r_1,c_1)$ & $(r_0+r_1, c_0+c_1)$ \\
\bottomrule
\end{tabular}
\end{table}

In these units, a roll changes only one coordinate; standing $\leftrightarrow$ lying shifts the
centre by 3 (1.5 cells) and lying $\to$ lying by 2; a cube step shifts the cubes' midpoint by 1;
\texttt{SWAP} shifts nothing. Every move of cost $\ge 1$ therefore moves the centre by at most 3 in
Manhattan distance, which gives
\begin{equation}
\mathrm{direct}(p, q) = \left\lceil \frac{|p_r - q_r| + |p_c - q_c|}{3} \right\rceil .
\end{equation}
This is the assignment's ``Manhattan distance divided by 2'' idea, tightened to 1.5 cells per roll.
It is also orientation-aware without an extra term: a lying block's centre sits half a cell off the
grid, so a lying block next to the goal still gets $h \ge 1$.

\textbf{Split switches.} A split switch teleports the block, so $\mathrm{direct}(p, \text{goal})$
alone could overestimate. For each split switch $q$ with entry $E_q$ (its \texttt{P} tile) and exit
$X_q$ (the midpoint of its destination cubes), the code computes $T[q]$, the best bound from $X_q$
to the goal, possibly through other switches, by Bellman--Ford style relaxation over the few
switches. The heuristic is
\begin{equation}
h(p) = \min\Big( \mathrm{direct}(p, \text{goal}),\ \min_q \big[\mathrm{direct}(p, E_q) + T[q]\big] \Big),
\end{equation}
treating the teleport as free, which can only lower the bound.

\subsubsection{Admissibility and consistency}
\label{sec:admissible}

The heuristic must be admissible with respect to our cost function, \texttt{move\_cost}.

\textbf{Admissible.} A path to the goal either never fires a split switch or fires one. In the
first case each move shifts the centre by at most 3, so at least $\mathrm{direct}(p,\text{goal})$
moves are needed, each costing at least 1. In the second case at least
$\mathrm{direct}(p, E_q)$ moves come before the teleport and at least $T[q]$ after it, by the same
argument from the exit. \texttt{SWAP} costs 0 but does not move the centre, and the fragile penalty
only increases the true cost. Hence $h(n) \le h^*(n)$ for every state.

\textbf{Consistent} ($h(s) \le c(s,s') + h(s')$ for every edge):
\begin{itemize}[leftmargin=*]
\item \emph{Roll or cube step:} $\mathrm{direct}$ satisfies the triangle inequality because
$\lceil (a+b)/3 \rceil \le \lceil a/3 \rceil + \lceil b/3 \rceil$, and one move shifts the centre by
at most 3, so $h(s) \le 1 + h(s') \le c + h(s')$ for any $c \ge 1$, including penalised moves.
\item \emph{SWAP:} the centre is unchanged, so $h(s) = h(s')$ and $c = 0$.
\item \emph{Teleport:} $h(s) \le \mathrm{direct}(p, E_q) + T[q] \le 1 + h(s')$, since $T[q]$ is
exactly the heuristic value at the exit.
\end{itemize}
A consistent heuristic is also admissible, and lets A* close each state once while remaining
optimal. The experiments agree: A* and UCS return the same cost on all 34 test stages.

\subsubsection{Limitations}

The bound ignores terrain, so it is loose on winding maps where the real path must detour, and it
prices every move at 1, so it underestimates when fragile tiles cannot be avoided. A tighter
admissible alternative would be a relaxed point-on-grid BFS distance computed once per stage.

% -------------------------------------------------------------------------
\subsection{Summary}

\begin{table}[H]
\centering
\caption{Theoretical comparison of the four algorithms}
\label{tab:theory}
\begin{tabularx}{\linewidth}{lXXXX}
\toprule
 & \textbf{BFS} & \textbf{DFS} & \textbf{UCS} & \textbf{A*} \\
\midrule
Frontier & FIFO queue & LIFO stack & heap on $g$ & heap on $g+h$ \\
Goal test & on generation & on pop & on pop & on pop \\
Complete & yes & yes (finite space) & yes & yes \\
Optimal for & number of actions & -- & \texttt{move\_cost} & \texttt{move\_cost} \\
Uses cost / heuristic & no / no & no / no & yes / no & yes / yes \\
Time & $O(b^d)$ & $O(b^m)$ & $O(b^{1+\lfloor C^*/\epsilon \rfloor})$ & depends on $h$ \\
\bottomrule
\end{tabularx}
\end{table}
